% TOP_PROBLEM: {"id": 10000047, "title": "Liouville property of infinite Ramanujan graphs", "category_id": 3, "category": {"id": 3, "name": "graph_theory", "display_name": "Graph Theory", "description": "Problems involving graphs, networks, and their properties.", "slug": "graph-theory", "order_index": 3, "created_at": "2026-07-31T15:26:25.671Z"}, "status": "open", "problem_number": "AMR-099-0047", "set_id": 13, "statusQualification": "Explicitly excludes degree two, where the infinite path is a Liouville Ramanujan graph and the literal unqualified assertion is false."}
% FIRST_POSTED: 2026-10-01
% ENTRY_KIND: statement_only
% UnsolvedMath ID 10000047; source catalogue code AMR-099-0047
% !TeX program = lualatex
% Definition and sources only; this file is not an LLM solution attempt.
\documentclass[11pt,a4paper]{article}
\usepackage[margin=25mm]{geometry}
\usepackage{fontspec}
\setmainfont{lmroman10-regular.otf}[BoldFont=lmroman10-bold.otf,
 ItalicFont=lmroman10-italic.otf,BoldItalicFont=lmroman10-bolditalic.otf]
\usepackage{amsmath,amssymb,mathtools}
\usepackage{unicode-math}
\setmathfont{latinmodern-math.otf}
\AtBeginDocument{\let\setminus\smallsetminus}
\usepackage{xurl,hyperref}
\hypersetup{colorlinks=true,urlcolor=blue}
\setcounter{secnumdepth}{0}
\setlength{\parindent}{0pt}
\setlength{\parskip}{5pt}
\setlength{\emergencystretch}{3em}
\begin{document}
\section*{Liouville property of infinite Ramanujan graphs}\label{10000047}
{\small UnsolvedMath ID 10000047 · AMR-099-0047 · Graph Theory · AMR Open Problem Lists}
\subsection{Definitions and mathematical statement}
Fix an integer $d\ge3$ and an infinite connected simple $d$-regular graph $G$. Its simple-random-walk operator acts on functions by
\[
(Pf)(v)=\frac1d\sum_{w\sim v}f(w).
\]
On $\ell^2(V(G))$ with counting measure this is a bounded self-adjoint operator. Call $G$ Ramanujan when
\[
\rho(G):=\|P\|_{\ell^2\to\ell^2}=\frac{2\sqrt{d-1}}{d},
\]
the spectral radius of the infinite $d$-regular tree. Equivalently, $\rho(G)=\limsup_{n\to\infty}P^n(v,v)^{1/n}$; the limsup handles possible periodicity.

A real function $f$ on $V(G)$ is harmonic if $Pf=f$. The graph is Liouville if every bounded harmonic real function is constant. Prove that no such infinite Ramanujan graph is Liouville: there must exist a nonconstant $f$ with $\sup_v|f(v)|<\infty$ and $Pf=f$.

No transitivity or group structure is assumed. The hypothesis is the sharp tree spectral-radius value, not merely positive expansion or a spectral radius strictly below one. The conclusion concerns bounded harmonic functions and therefore differs from existence of an arbitrary unbounded harmonic function.

The degree restriction removes a necessary degenerate exception to the source's shorthand: the two-way infinite path is $2$-regular, has spectral radius one equal to its tree value, and is Liouville. The intended nonamenable Ramanujan question is precisely the range $d\ge3$, for which the displayed tree spectral radius is strictly less than one.
\subsection{Short English statement}
Must every infinite Ramanujan graph of degree at least three admit a nonconstant bounded harmonic function?
\subsection{Sources}
\begin{itemize}
\item[S1] ULAM.AI and the UnsolvedMath Contributors, supplied problems.json, record 10000047 (AMR-099-0047), AMR Open Problem Lists. Catalogue identity and source transcription; CC BY 4.0.
\url{https://huggingface.co/datasets/ulamai/UnsolvedMath}.
\item[S2] Benjamini - Coarse Geometry and Randomness (Saint-Flour notes), Conjecture 13.1, PDF page 107. Original source indexed by AMR-099-0047.
\url{https://www.wisdom.weizmann.ac.il/~itai/stflouraug24.pdf#page=107}.
\item[S3] I. Benjamini, Coarse Geometry and Randomness, primary Saint-Flour notes; the numbered question and its surrounding definitions.
\url{https://arquivo.pt/noFrame/replay/20201231041548id_/http://www.wisdom.weizmann.ac.il/~itai/stflouraug24.pdf}.
\end{itemize}
\end{document}
